% ============================================================
%  Lista de Exercícios — Template | PratiqueJá
%  Engine : XeLaTeX
%  Layout : \listheader em todas as páginas → overhead uniforme
%           3 linhas × 2 colunas = 6 exercícios por página
%           3 páginas → 18 exercícios
%
%  Calibração de ht:
%    compile com \showsolucaofalse
%    3 páginas = correto, 4 páginas = ht grande demais
%
%  Versão Grátis:  comente a linha \showsolucaotrue
%  Versão Premium: mantenha  a linha \showsolucaotrue
% ============================================================
\documentclass[12pt]{article}
\usepackage{../pratiqueja}
\usepackage{tabularray}

% ── 1. PERSONALIZE ESTAS LINHAS ──────────────────────────────
\newcommand{\lstTitulo}{Área e Perímetro}
\newcommand{\lstSubtitulo}{Nível III}
\newcommand{\lstCategoria}{Matemática · Básico}
\newcommand{\lstInstrucao}{%
  Resolva as questões abaixo mostrando o desenvolvimento completo.}

% ── 2. CHAVE GRÁTIS / PREMIUM ────────────────────────────────
\newif\ifshowsolucao
\showsolucaotrue

\setsubject{\lstTitulo\ · \lstSubtitulo}

% ============================================================
%  MACROS DE EXERCÍCIO
% ============================================================

\newcommand{\numex}[1]{%
  \begin{tcolorbox}[
    enhanced, arc=3pt, boxrule=0pt,
    colback=rulebg, colframe=rulebg,
    left=4pt, right=4pt, top=1pt, bottom=1pt,
    nobeforeafter, on line,
  ]{\bfseries\small\color{darkblue}#1}\end{tcolorbox}%
  \hspace{4pt}%
}

\newcommand{\altline}[2]{%
  \def\altph{---}%
  \def\altarg{#2}%
  \ifx\altarg\altph\else
    \noindent{\bfseries\color{darkblue}\small(#1)}~{\small#2}\par
  \fi}

\newcommand{\alts}[5]{%
  \altline{A}{#1}%
  \altline{B}{#2}%
  \altline{C}{#3}%
  \altline{D}{#4}%
  \altline{E}{#5}}

\newlength{\cellheight}
\setlength{\cellheight}{7.8cm}

\newcommand{\ex}[7]{%
  \parbox[t][\cellheight][t]{\linewidth}{%
    \setlength{\parskip}{6pt}%
    \noindent\numex{#1}{\small\color{bodytext}#2}\par\vfill
    \alts{#3}{#4}{#5}{#6}{#7}}}

\newcommand{\exd}[2]{%
  \noindent\numex{#1}{\small\color{bodytext}#2}%
}

% ============================================================
%  MACROS DE GABARITO
% ============================================================

% card combinado gabarito: [01 · B]
\newcommand{\resheader}[2]{%
  \begin{tcolorbox}[enhanced,arc=3pt,boxrule=0pt,colback=rulebg,colframe=rulebg,
    left=4pt,right=4pt,top=1pt,bottom=1pt,nobeforeafter,on line]
    {\bfseries\small\color{darkblue}#1%
     \hspace{5pt}{\color{mutedtext}\tiny\textbullet}\hspace{5pt}%
     \color{vegreen}#2}%
  \end{tcolorbox}\hspace{4pt}}

\newcommand{\res}[3]{%
  \noindent\resheader{#1}{#2}\hspace{0.3cm}%
  {\fontsize{10}{12}\selectfont\setlength{\lineskip}{6pt}\setlength{\lineskiplimit}{2pt}\setlength{\jot}{8pt}\color{bodytext}#3\par}%
  \vspace{0.7cm}}

\newcommand{\resd}[2]{%
  \noindent\numex{#1}\hspace{0.3cm}%
  {\fontsize{10}{12}\selectfont\setlength{\lineskip}{6pt}\setlength{\lineskiplimit}{2pt}\setlength{\jot}{8pt}\color{bodytext}#2\par}%
  \vspace{0.7cm}}

% ============================================================
\begin{document}
\pagenumbering{arabic}
\setstretch{1.2}
\color{bodytext}

% ════════════════════════════════════════════════════════════
%  PÁGINA 1 — Exercícios 01–06  (3 linhas × ht=7.8cm)
% ════════════════════════════════════════════════════════════
\listheader{\lstTitulo}{\lstSubtitulo}{\lstCategoria}{\lstInstrucao}

\begin{tblr}{
  colspec  = {p{0.455\textwidth}|[1.5pt,sepcolor]p{0.455\textwidth}},
  columns  = {colsep=0pt, leftsep=7pt, rightsep=7pt},
  rows     = {ht=7.8cm, valign=t, rowsep=0pt},
}
\ex{01}{Enunciado do exercício 01 com um texto um pouco mais longo para testar.}{$10\ \text{cm}^2$}{$20\ \text{cm}^2$}{$30\ \text{cm}^2$}{$40\ \text{cm}^2$}{---} &
\ex{02}{Enunciado curto.}{A}{B}{C}{D}{E} \\
\ex{03}{Enunciado do exercício 03.}{$5\ \text{cm}$}{$10\ \text{cm}$}{$15\ \text{cm}$}{$20\ \text{cm}$}{---} &
\ex{04}{Enunciado do exercício 04 também com texto mais longo para ver o alinhamento das alternativas no fundo da célula.}{A}{B}{C}{D}{---} \\
\exd{05}{Enunciado do exercício 05 discursivo.} &
\exd{06}{Enunciado do exercício 06 discursivo.} \\
\end{tblr}

% ════════════════════════════════════════════════════════════
%  PÁGINA 2 — Exercícios 07–12  (3 linhas × ht=7.8cm)
% ════════════════════════════════════════════════════════════
\newpage
\listheader{\lstTitulo}{\lstSubtitulo}{\lstCategoria}{\lstInstrucao}

\begin{tblr}{
  colspec  = {p{0.455\textwidth}|[1.5pt,sepcolor]p{0.455\textwidth}},
  columns  = {colsep=0pt, leftsep=7pt, rightsep=7pt},
  rows     = {ht=7.8cm, valign=t, rowsep=0pt},
}
\exd{07}{Enunciado do exercício 07.} &
\exd{08}{Enunciado do exercício 08.} \\
\exd{09}{Enunciado do exercício 09.} &
\exd{10}{Enunciado do exercício 10.} \\
\exd{11}{Enunciado do exercício 11.} &
\exd{12}{Enunciado do exercício 12.} \\
\end{tblr}

% ════════════════════════════════════════════════════════════
%  PÁGINA 3 — Exercícios 13–18  (3 linhas × ht=7.8cm)
% ════════════════════════════════════════════════════════════
\newpage
\listheader{\lstTitulo}{\lstSubtitulo}{\lstCategoria}{\lstInstrucao}

\begin{tblr}{
  colspec  = {p{0.455\textwidth}|[1.5pt,sepcolor]p{0.455\textwidth}},
  columns  = {colsep=0pt, leftsep=7pt, rightsep=7pt},
  rows     = {ht=7.8cm, valign=t, rowsep=0pt},
}
\exd{13}{Enunciado do exercício 13.} &
\exd{14}{Enunciado do exercício 14.} \\
\exd{15}{Enunciado do exercício 15.} &
\exd{16}{Enunciado do exercício 16.} \\
\exd{17}{Enunciado do exercício 17.} &
\exd{18}{Enunciado do exercício 18.} \\
\end{tblr}

% ════════════════════════════════════════════════════════════
%  GABARITO (somente Premium)
% ════════════════════════════════════════════════════════════
\ifshowsolucao
\clearpage

\heroheader{Gabarito — \lstTitulo}%
           {\lstSubtitulo\ · Resolução comentada}%
           {\lstCategoria}
\vspace{4pt}

\DefTblrTemplate{caption}{default}{}
\DefTblrTemplate{conthead}{default}{}
\DefTblrTemplate{conthead-text}{default}{}
\DefTblrTemplate{contfoot}{default}{}
\DefTblrTemplate{contfoot-text}{default}{}
\begin{longtblr}[label=none]{
  colspec  = {p{0.455\textwidth}|[1.5pt,sepcolor]p{0.455\textwidth}},
  columns  = {colsep=0pt, leftsep=7pt, rightsep=7pt},
  rows     = {ht=3.5cm, valign=t, rowsep=0pt},
}
\res{01}{B}{%
  Área do retângulo: base $\times$ altura\\
  $A = \dfrac{b \cdot h}{2} = \dfrac{8 \cdot 5}{2} = \dfrac{40}{2} = \mathbf{20}\ \text{cm}^2$} &
\res{02}{C}{%
  Perímetro: soma de todos os lados\\
  $P = 2(b + h) = 2(6 + 4) = 2 \cdot 10 = \mathbf{20}\ \text{cm}$} \\
\res{03}{A}{%
  Área do triângulo:\\
  $A = \dfrac{b \cdot h}{2} = \dfrac{12 \cdot 7}{2} = \dfrac{84}{2} = \mathbf{42}\ \text{cm}^2$} &
\res{04}{D}{%
  Circunferência: $C = 2\pi r$\\
  $C = 2 \cdot \pi \cdot 5 = 10\pi \approx \mathbf{31{,}4}\ \text{cm}$} \\
\res{05}{E}{%
  Área do círculo:\\
  $A = \pi r^2 = \pi \cdot 3^2 = 9\pi \approx \mathbf{28{,}3}\ \text{cm}^2$} &
\res{06}{B}{%
  Trapézio — bases $B=10$, $b=6$, altura $h=4$:\\
  $A = \dfrac{(B+b)\cdot h}{2} = \dfrac{(10+6)\cdot 4}{2} = \dfrac{64}{2} = \mathbf{32}\ \text{cm}^2$} \\
\resd{07}{%
  Losango — diagonais $d_1=8$, $d_2=6$:\\
  $A = \dfrac{d_1 \cdot d_2}{2} = \dfrac{8 \cdot 6}{2} = \dfrac{48}{2} = \mathbf{24}\ \text{cm}^2$} &
\resd{08}{%
  Quadrado de lado $\ell = 7\ \text{cm}$:\\
  $A = \ell^2 = 7^2 = \mathbf{49}\ \text{cm}^2$\\
  $P = 4\ell = 4 \cdot 7 = \mathbf{28}\ \text{cm}$} \\
\resd{09}{Resolução 09} &
\resd{10}{Resolução 10} \\
\resd{11}{Resolução 11} &
\resd{12}{Resolução 12} \\
\resd{13}{Resolução 13} &
\resd{14}{Resolução 14} \\
\resd{15}{Resolução 15} &
\resd{16}{Resolução 16} \\
\resd{17}{Resolução 17} &
\resd{18}{Resolução 18} \\
\end{longtblr}
\fi

\end{document}
